\documentclass[11pt]{amsart}

\usepackage{booktabs}
\usepackage[hidelinks]{hyperref}

\newtheorem{theorem}{Theorem}
\newtheorem{lemma}{Lemma}[section]

\title{Skolem pairings with separated pair sums}
\author{Gleb Golovkov}
\subjclass[2020]{Primary 05A18, Secondary 05B30}
\keywords{Skolem sequence, pair partition, separated sums, interval certification}

\begin{document}

\begin{abstract}
For \(n\ge2\), the integers \(1,\ldots,2n\) admit a pairing with
differences \(1,\ldots,n\) and pair sums separated by at least two
if and only if \(n\ge5\) and \(n\equiv0,1\pmod4\).
Parity and the case \(n=4\) give necessity.
Sufficiency uses linear-periodic block constructions in 18 classes
modulo 36, machine-certified interval boundaries, and a finite
search for 74 small orders.
\end{abstract}

\maketitle

\section{Introduction}

Write \([m]=\{1,\ldots,m\}\). A \emph{Skolem pairing of order \(n\)}
is a partition of \([2n]\) into pairs \((a_k,b_k)\) with
\(b_k-a_k=k\) for \(1\le k\le n\). We impose a separation condition
on the ordinary integer sums \(s_k=a_k+b_k\).

\begin{theorem}\label{thm:main}
Let \(n\ge2\) be an integer. There is a Skolem pairing of order \(n\) with
\[
  |s_i-s_j|\ge2\qquad(1\le i<j\le n)
\]
if and only if \(n\ge5\) and \(n\equiv0,1\pmod4\).
\end{theorem}

A Skolem pairing gives a Skolem sequence by placing the two
copies of \(k\) at positions \(a_k,b_k\) in a word of length
\(2n\). Conversely, the positions of each symbol in a Skolem
sequence recover its pair. In this language the theorem
requires the midpoints of any two position pairs to be at
distance at least one.

For \(n=5\), the pairs in increasing order of difference are
\[
  (8,9),\quad (1,3),\quad (4,7),\quad (2,6),\quad (5,10).
\]
The sorted sums are \((4,8,11,15,17)\). Their successive
gaps are \((4,3,4,2)\).

Skolem proved that the difference condition alone is possible
exactly when \(n\equiv0,1\pmod4\) \cite{skolem}. A strong Skolem
starter requires the pair sums to have distinct nonzero residues
modulo \(2n+1\). The example modulo 11 in the preprint of
Ogandzhanyants, Kondratieva, and Shalaby has consecutive ordinary
sums 15 and 16 \cite[Example~1.6]{strong}.

The repository address for the computational supplement is\newline
\url{https://github.com/gleb-ml/spaced-skolem-partitions}.

\section{Necessity}

There are \(n\) odd labels in \([2n]\). Each of the
\(\lceil n/2\rceil\) pairs of odd difference uses exactly one odd
label. Every remaining pair has even difference, so it uses either
two odd labels or two even labels. Consequently
\(n-\lceil n/2\rceil=\lfloor n/2\rfloor\) is even. This holds
exactly when \(n\equiv0,1\pmod4\).

It remains to exclude \(n=4\). The reflection \(x\mapsto9-x\)
preserves differences and changes each sum \(s\) into \(18-s\).
Up to reflection, the pair of difference 4 is \((1,5)\) or
\((2,6)\). We list the possible completions by choosing the
pairs of differences 3 and 2 in that order.

If the pair of difference 4 is \((1,5)\), the choices for
difference 3 are \((3,6)\) and \((4,7)\). After \((3,6)\),
only \((2,4)\) has difference 2 among the free labels, leaving
\((7,8)\). After \((4,7)\), only \((6,8)\) has difference 2,
leaving \((2,3)\).

If the pair of difference 4 is \((2,6)\), the choices for
difference 3 are \((1,4)\), \((4,7)\), and \((5,8)\). After \((1,4)\) the
pairs of difference 2 are \((3,5)\) and \((5,7)\), which leave
\((7,8)\) and \((3,8)\), respectively. After \((4,7)\) the
pairs of difference 2 are \((1,3)\) and \((3,5)\), which leave
\((5,8)\) and \((1,8)\), respectively. After \((5,8)\) the
only pair of difference 2 is \((1,3)\), leaving \((4,7)\).
Thus the only completion in this branch uses \((1,4)\), \((3,5)\),
and \((7,8)\).

The three resulting lists of sums in increasing order are
\[
  (6,6,9,15),\qquad (5,6,11,14),\qquad (5,8,8,15).
\]
Each list contains a repeated sum or consecutive sums. Reflection
preserves this obstruction.

\section{Linear interval certification}

\subsection{Block parameters}

Put \(n=12t+r\) and let \(c\) be the residue of \(t\) modulo 3,
with \(0\le r<12\) and \(0\le c<3\). In each class \((r,c)\),
a block consists of the pairs
\begin{equation}\label{eq:block}
  (At+a+pj,\ Bt+b+qj),\qquad j\in\mathbb Z,\quad 0\le j<wt+\ell.
\end{equation}
The ten shapes \((A,p,B,q,w)\) in Table~\ref{tab:shapes} are
common to all classes. Each class specifies a triple
\((a,b,\ell)\) for each shape and a list of individual pairs
\((At+a,Bt+b)\), given by quadruples \((A,a,B,b)\).
For each block, its left ends, right ends, differences, and sums
are arithmetic progressions with signed steps \(p,q,q-p,p+q\),
respectively. An individual pair contributes a singleton to each
of these lists.

\begin{table}[htbp]
\centering
\caption{The common block shapes and their signed steps.}
\label{tab:shapes}
\begin{tabular}{@{}rrrrrrrr@{}}
\toprule
Block & \(A\) & \(p\) & \(B\) & \(q\) & \(w\) & \(q-p\) & \(p+q\)\\
\midrule
1 & 10 & -4 & 20 & -2 & 1 & 2 & -6\\
2 & 12 & -4 & 20 & -2 & 1 & 2 & -6\\
3 & 18 & -4 & 24 & -2 & 1 & 2 & -6\\
4 & 8 & -4 & 14 & 2 & 1 & 6 & -2\\
5 & 10 & -4 & 10 & 2 & 2 & 6 & -2\\
6 & 2 & -2 & 2 & 4 & 1 & 6 & 2\\
7 & 4 & -2 & 4 & 4 & 2 & 6 & 2\\
8 & 12 & 2 & 14 & 4 & 1 & 2 & 6\\
9 & 16 & 2 & 20 & 4 & 1 & 2 & 6\\
10 & 20 & 2 & 20 & 4 & 1 & 2 & 6\\
\bottomrule
\end{tabular}
\end{table}

\subsection{Interval lemma}

For an integer interval \([L,R]\) we mean the set of integers
from \(L\) to \(R\), inclusive. A list of progressions or
intervals is treated as a family with indices, so identical
entries contribute distinct copies.

\begin{lemma}[Linear interval certification]\label{lem:certify}
Let \(t=t_*+12u\) for \(u\in\mathbb Z_{\ge0}\), with \(t_*\)
a fixed integer. Consider a finite family of progressions
\[
  \{\alpha t+\beta+ej:j\in\mathbb Z,\ 0\le j<wt+\ell\}
\]
with integer coefficients. Each either has exactly one term
for all \(t\), or has \(w>0\), \(d=|e|>0\) dividing 12, and
\(wt_*+\ell\ge12/d\). For each residue \(z\) modulo 12, take
the intervals obtained from the progressions in the coordinate
\(v=(x-z)/12\) and order them by their left endpoints at \(u=0\).
Use this order for all subsequent comparisons.

\begin{enumerate}
\item If the gaps between consecutive intervals are nonnegative
at \(u=0\) and each gap at \(u=1\) is at least its value at
\(u=0\), the progressions are pairwise disjoint for every
\(u\ge0\).
\item Suppose \(M(t)\) is affine in \(t\) with integer coefficients
and \(M(t)\ge12\) for all \(u\ge0\). If each residue has a
nonempty list and the gaps between consecutive intervals and
both gaps at the ends of \([M(t)]\) are zero at \(u=0,1\),
the progressions partition \([M(t)]\) for every \(u\ge0\).
\end{enumerate}
\end{lemma}

\begin{proof}
Write \(m=wt+\ell\). For a progression with \(e\ne0\), its two
extreme terms are \(\alpha t+\beta\) and
\(\alpha t+\beta+e(m-1)\). Their order is fixed because
\(m\ge12/d\ge1\) and the sign of \(e\) is fixed. Thus the
ordered extremes are
\[
  \operatorname{lo}(u)=L_0+12\lambda u,\qquad
  \operatorname{hi}(u)=R_0+12\rho u
\]
for integers \(L_0,R_0,\lambda,\rho\).
The progression lies in a single residue modulo \(d\),
which is independent of \(u\). In each run of \(12/d\)
consecutive terms it visits every compatible residue modulo 12
exactly once. Hence a compatible residue occurs at \(u=0\)
and at all later values of \(u\). For such a residue \(z\),
the interval in the coordinate \(v\) has endpoints
\begin{align*}
  L(u)&=\left\lceil\frac{\operatorname{lo}(u)-z}{12}\right\rceil
       =\left\lceil\frac{L_0-z}{12}\right\rceil+\lambda u,\\
  R(u)&=\left\lfloor\frac{\operatorname{hi}(u)-z}{12}\right\rfloor
       =\left\lfloor\frac{R_0-z}{12}\right\rfloor+\rho u.
\end{align*}
A singleton belongs to a fixed residue modulo 12 and gives an
interval whose two endpoints are the same affine function.
Thus the list of intervals for each \(z\) stays fixed as a list
of indexed affine endpoints.

For consecutive intervals in the chosen order, the number of
integers strictly between them, when they are disjoint, is
measured by the signed gap
\[
  g(u)=L_{\mathrm{next}}(u)-R_{\mathrm{prev}}(u)-1
      =g(0)+u\bigl(g(1)-g(0)\bigr).
\]
The inequalities in part~(1) give \(g(u)\ge0\), so this order
remains disjoint. Distinct residues are disjoint as sets of
integers, which proves part~(1).

For part~(2), the target interval in residue \(z\) is
\[
  \left[\left\lceil\frac{1-z}{12}\right\rceil,
         \left\lfloor\frac{M(t)-z}{12}\right\rfloor\right].
\]
The end gaps are the first left endpoint minus the target left
endpoint and the target right endpoint minus the last right
endpoint. They are affine in \(u\), since \(M(t)\) changes by a
multiple of 12 when \(t\) increases by 12. Each end gap and each
internal gap vanishes identically if it is zero at \(u=0,1\).
The intervals therefore partition the target in every residue.
\end{proof}

\subsection{Application to ends, differences, and sums}

The absolute values of the four steps in each row of
Table~\ref{tab:shapes} belong to \(\{2,4,6\}\), and each
\(w\) is positive. The hypothesis on the number of terms is
checked at \(t=t_*\). For the ends, we apply part~(2) of the
lemma to the combined family of left and right progressions,
with target \([2n]\). For the differences, we apply part~(2)
with target \([n]\). This also forces \(a_k<b_k\) in every
pair. The difference can then index the pair uniquely.

To test separation, take each progression \(S\) of sums and its
translate \(S+1\), retaining their indices. The signed sum step
is even and nonzero in every block, so the elements within
\(S\) are distinct, and \(S\) and \(S+1\) have opposite
parity. For a singleton, its sum and its translate also have
opposite parity. Hence a progression and its own translate
do not create a conflict in any residue modulo 12.
Two equal sums from distinct pairs appear in the same residue
in the unshifted family. If two original sums differ by one,
the larger equals an element of the translated family.
Disjointness of the combined family therefore certifies
\(|s_i-s_j|\ge2\). We obtain this disjointness by applying
part~(1) of Lemma~\ref{lem:certify}.

\section{Completion of the proof, verification, and open problems}

\subsection{The certificate}

The 18 classes and their thresholds are listed in
Table~\ref{tab:thresholds}. The pairs \((r,c)\) correspond to the
residues \(r+12c\) modulo 36. For a row with threshold \(t_0\),
every \(t\ge t_0\) in the class is covered by
\[
  t=t_*+12u,\qquad
  t_*\in\{t_0,t_0+3,t_0+6,t_0+9\},\quad
  u\in\mathbb Z_{\ge0}.
\]
The full 18 tables of triples and quadruples are part of the
computational supplement. They are embedded in \texttt{proof.py}
and printed in plain text in \texttt{tables.txt}. The latter is
generated by \texttt{python3 proof.py --tables}.

\begin{table}[htbp]
\centering
\caption{Certified thresholds. The final column counts individual
pairs in addition to the ten blocks.}
\label{tab:thresholds}
\begin{tabular}{@{}rrrrr@{}}
\toprule
\(r\) & \(c\) & \(t_0\) & \(n_0=12t_0+r\) & Individual pairs\\
\midrule
0 & 0 & 9 & 108 & 20\\
0 & 1 & 10 & 120 & 28\\
0 & 2 & 11 & 132 & 30\\
1 & 0 & 15 & 181 & 30\\
1 & 1 & 10 & 121 & 31\\
1 & 2 & 17 & 205 & 30\\
4 & 0 & 9 & 112 & 28\\
4 & 1 & 13 & 160 & 54\\
4 & 2 & 17 & 208 & 71\\
5 & 0 & 15 & 185 & 49\\
5 & 1 & 19 & 233 & 44\\
5 & 2 & 14 & 173 & 33\\
8 & 0 & 15 & 188 & 52\\
8 & 1 & 10 & 128 & 34\\
8 & 2 & 17 & 212 & 36\\
9 & 0 & 15 & 189 & 53\\
9 & 1 & 16 & 201 & 48\\
9 & 2 & 17 & 213 & 48\\
\bottomrule
\end{tabular}
\end{table}

The routine \texttt{certify} forms the intervals at \(t_*\)
and \(t_*+12\), using the order of the intervals at \(t_*\)
for both evaluations. Across the 72 cases it checks
\(30{,}076\) internal or end gaps, each at both parameter
values. For the ends and differences, the two evaluations of
each gap are zero. For the combined family of sums and shifted
sums, the first evaluation is nonnegative and the second is at
least the first. These checks have also been reproduced by
computing the affine interval endpoints directly from the first
and last indices in each residue class.

\begin{proof}[Proof of sufficiency in Theorem~\ref{thm:main}]
For \(n\) at or above the threshold of its class, the
certificate and Lemma~\ref{lem:certify} give a partition with
the three required properties. Every admissible \(n\) below its
own class threshold is less than \(233\). There are 74 such
orders. For each, the script finds a pairing by finite search
and checks the sorted endpoint list, sorted difference list,
and successive gaps of the sorted sums by direct evaluation.
Together the two parts cover every \(n\ge5\) with
\(n\equiv0,1\pmod4\).
\end{proof}

The script additionally runs \texttt{verify} for every
admissible \(n\) from 5 through 600. It checks 298 pairings,
including 224 obtained directly from the block formulas.

\subsection{Enumeration at small orders}

Table~\ref{tab:counts} counts each pairing once and counts
distinct reflections separately. The enumeration assigns
differences in decreasing order. It tries each free left end
and its determined right end, discarding branches with a
reused label or a new sum within one of an earlier sum.
Each valid pairing determines one branch.

\begin{table}[htbp]
\centering
\caption{Number of Skolem pairings with separated pair sums.}
\label{tab:counts}
\begin{tabular}{@{}lrrrrrrrrr@{}}
\toprule
Order \(n\)  & 2 & 3 & 4 & 5 & 6 & 7 & 8 & 9 & 12\\
\midrule
Pairings  & 0 & 0 & 0 & 2 & 0 & 0 & 6 & 30 & 1300\\
\bottomrule
\end{tabular}
\end{table}

\subsection{Open problems}

\begin{enumerate}
\item For \(d\ge1\), replace the prescribed differences by
\(d,d+1,\ldots,d+n-1\), as in the Langford problem
\cite{langford,simpson}. For ordinary pairings Simpson's
criterion is \(n\ge2d-1\), with \(n\equiv0,1\pmod4\) for
odd \(d\) and \(n\equiv0,3\pmod4\) for even \(d\).
At \(n=2d-1\), the differences sum to \(n^2\), forcing the
left ends to be \([n]\) and the right ends its complement.
Separation in \([n+2,3n]\) would give the sums one parity,
but each sum has the parity of its difference and both
parities occur for \(d\ge2\).
Under Simpson's parity conditions, does \(n\ge2d+3\)
suffice for separation apart from \((d,n)=(2,7)\)?
The latter case has zero separated pairings by exhaustive
enumeration.

\item Can one construct the separated pairings with fewer
congruence classes, or specify the individual pairs by a
uniform rule?

\item Along \(n\equiv0,1\pmod4\), what is the asymptotic
number of separated pairings? In particular, does their
proportion among all Skolem pairings tend to zero?
\end{enumerate}

\subsection*{Acknowledgments}
AI-assisted tools were used to search for block parameters and
to organize the manuscript.
The mathematical arguments and the verification script were
checked independently, including by a separate implementation
of the interval checks.

\begin{thebibliography}{9}

\bibitem{langford}
C.~D. Langford, \emph{Problem},
Math. Gaz. \textbf{42} (1958), no.~341, 228.
\href{https://doi.org/10.2307/3610395}{doi:10.2307/3610395}.

\bibitem{strong}
O.~Ogandzhanyants, M.~Kondratieva, and N.~Shalaby,
\emph{Strong Skolem starters},
J. Combin. Des. \textbf{27} (2019), no.~1, 5--21.
\href{https://doi.org/10.1002/jcd.21633}{doi:10.1002/jcd.21633}.
Preprint \href{https://arxiv.org/abs/1805.05293}{arXiv:1805.05293}.

\bibitem{simpson}
J.~E. Simpson, \emph{Langford sequences: perfect and hooked},
Discrete Math. \textbf{44} (1983), no.~1, 97--104.
\href{https://doi.org/10.1016/0012-365X(83)90008-0}
     {doi:10.1016/0012-365X(83)90008-0}.

\bibitem{skolem}
T.~Skolem, \emph{On certain distributions of integers in pairs
with given differences}, Math. Scand. \textbf{5} (1957), 57--68.
\href{https://doi.org/10.7146/math.scand.a-10490}
     {doi:10.7146/math.scand.a-10490}.

\end{thebibliography}

\end{document}
